\chapter{日本企業データ構築と変数定義：四要素モデルの実証的写像}
\label{chap:data_construction}

\section{本章の目的と基本方針}

本章では，第4章で導出した倒産距離 $Z_{\text{lin}}$ の構造
（収益性 $\mu$，市場ショック感応度 $\sigma$，支援能力 $S$，倒産閾値の代理変数 $D$）
を，日本の新電力企業の財務データ・市場データに写像するための
\textbf{データ構築・変数定義・前処理} を体系的に示す。

本研究の実証分析は，以下の三段階で進められる。

\begin{enumerate}
  \item 財務データ・市場データ・倒産データを統合し，
        四要素モデルの構成要素を観測可能な形に整形した
        月次パネルデータ（\texttt{financial\_data.csv}）を構築する。
  \item 構築したデータを用いて企業ごとの倒産距離 $DD_{i,t}$ を算出し，
        新電力企業の倒産リスク構造をミクロレベルで把握する。
  \item 7社全体の記述統計を求め，業界全体の財務特性や変数分布をマクロレベルで整理し，
        第9章の Ordered Logit 推定に必要なデータの性質を確認する。
\end{enumerate}

この順番は，倒産距離モデルの構造（平均 $-$ 閾値）／不確実性に沿って，
ミクロ分析からマクロ分析へと段階的に進むためのものである。
すなわち，月次パネルデータは理論モデルの四要素
（収益性 $\mu$，市場ショック $\sigma$，支援能力 $S$，倒産閾値の代理変数 $D$）
を観測可能な形に写像した基礎データであり，
企業別倒産距離 $DD$ はミクロの倒産リスク指標，
記述統計は業界全体の構造的特徴を把握するためのマクロ指標として位置づけられる。

---

\section{データ収集の全体像}

本研究で用いるデータは以下の三種類で構成される。

\begin{enumerate}
\item \textbf{財務データ}：官報に掲載された新電力企業の決算公告（年次）
\item \textbf{市場データ}：JEPX スポット市場価格（30分値・日次・月次）
\item \textbf{倒産データ}：破産・民事再生・登録取消・事業撤退の発生月
\end{enumerate}

財務データが取得可能であった新電力企業は 7 社であり，
分析期間は 2016 年 10 月から 2025 年 5 月までである。
新電力企業は非上場である場合が多く，有価証券報告書が入手できないため，
官報公告は財務情報を得るうえで最も信頼性の高い公開情報源である。

倒産段階（default\_stage）は，観測可能な退出イベントに基づき，
以下の 3 種類のイベントを基準として定義する。

\begin{itemize}
\item 法的倒産（破産・民事再生）
\item 小売電気事業の登録取消（経済産業省）
\item 事業撤退（企業発表・報道）
\end{itemize}

これらのイベントの発生日を月次パネルデータに割り当てることで，
企業 $i$ の月 $t$ における倒産段階 $\text{default\_stage}_{i,t}$ を構築した。

---

\section{財務データの処理}

\subsection{抽出項目}

官報掲載の決算公告から以下の項目を手作業で抽出した。

\begin{itemize}
\item 資産合計，負債合計，純資産
\item 流動資産，固定資産
\item 流動負債，固定負債
\item 資本金，利益剰余金
\end{itemize}

これらは倒産距離 DD の構成要素である
収益性 $\mu$，倒産閾値の代理変数 $D$ を構築するための基礎データである。

\subsection{年次データの月次化}

財務データは年次であるため，以下の手順で月次へ変換した。

\begin{enumerate}
\item 年次データを観測月に割り当てる（例：2020年決算 → 2020年3月）
\item 年次間を線形補間し，月次データを生成する
\end{enumerate}

線形補間は，年度間の財務指標の変化を滑らかに再現し，
外れ値の影響を抑えつつ趨勢を自然に反映できるため，
実証研究において広く用いられる方法である。

---

\section{市場データ：スポット価格の推移}

本研究で使用したスポット市場データ（システムプライス）は，
日本卸電力取引所（JEPX）が公開するスポット市場（Day Ahead Market）の価格情報に基づく。
市場価格の急変動は新電力企業の収益に直接影響するため，
本研究では日次変化率とその標準偏差（ボラティリティ）を計算し，
市場ショック感応度 $\sigma$ の実証的 proxy として用いる。

スポット価格の変化率は



\[
r_t = \frac{P_t - P_{t-1}}{P_{t-1}}
\]



と定義し，その標準偏差を市場ショックの大きさとして定量化する。
変化率の標準偏差は価格水準の違いを取り除き，
純粋な変動性のみを抽出できるため，市場ショックの強度を測定する指標として適切である。

---

\section{理論モデルの四要素の実証的写像}

本節では，理論モデルで導出した四要素
（収益性 $\mu$，市場ショック感応度 $\sigma$，支援能力 $S$，財務健全性 $D$）
を，実際の企業データに基づく観測可能な指標へと結び付ける。
これら四要素は倒産距離 $Z_{\text{lin}}$ の構造を形成する中心的な構成要素であり，
理論的には連続的な潜在変数として定義されているが，
実証分析では財務諸表・市場データ・企業属性情報などの具体的なデータを用いて
その意味内容を定量化する必要がある。

本節の目的は，各要素がどのような実証的 proxy によって測定され，
どのような経済的根拠に基づいて選択されるのかを体系的に示すことである。
これにより，理論モデルで構築した倒産リスクの構造が，
実証データの世界においてどのように再現されるかが明確になり，
後続の Ordered Logit 推定における線形部分 $X\beta$ との対応関係が
一貫した形で理解できるようになる。

% ------------------------------
% S
% ------------------------------

\subsection{支援能力 $S$ の定義と採用根拠}

支援能力 $S$ は，企業が財務的困難に直面した際に，
親会社や主要株主からどの程度の支援を受けられるかを表す指標である。
理論モデルでは，外部支援の余力が倒産閾値を押し下げ，
倒産距離を拡大させる方向に作用することが明確に示されている。

しかし，本研究で対象とする新電力企業では，
親会社支援が財務指標に明示的に反映されるケースが極めて少なく，
支援の有無や規模を定量的に測定できる情報が存在しない。
また，親会社が存在していても，
支援が制度的に保証されているわけではなく，
財務諸表や市場データから支援能力を推定することは困難である。

このため，本研究では支援能力を実証的に扱う際に
「観測可能な支援効果は存在しない」と解釈し，
$S_i = 0$ を採用する。
この設定により，倒産距離の構造は
収益性 $\mu$，市場ショック感応度 $\sigma$，財務健全性 $D$ の三要素によって決定され，
支援能力が倒産リスクを左右する可能性は理論的には保持しつつも，
実証モデルでは観測可能な情報に基づく簡潔な構造として整理される。

% ------------------------------
% D
% ------------------------------

\subsection{倒産閾値の代理変数 $D$（負債比率）の定義と採用根拠}

倒産閾値の代理変数 $D$ は，企業が保有する資産のうち，
どれだけが負債によって賄われているかを示す財務健全性の基本指標であり，



\[
D_{i,t}=\frac{\text{Liabilities}_{i,t}}{\text{Assets}_{i,t}}
\]



と定義される。

理論モデルでは倒産閾値 $\mathrm{Thr}$ は負債の名目価値に依存するが，
非上場企業では企業価値 $V$ や負債の市場価値を観測できないため，
倒産閾値を直接推定することができない。
このため，本研究では観測可能な財務指標である負債比率